\chapter{Chapter06 — Random Functions}
%%% generated by the blueprint pipeline from FABL.Chapter06.Pseudorandomness.RandomFunctions
\section{Overview}

Declarations of the Lean module \texttt{FABL.Chapter06.Pseudorandomness.RandomFunctions}.

\section{Declarations}

\begin{definition}[Boolean-valued function on the sign cube]
\lean{FABL.ZeroOneFunction}
\statementsource{Analysis-of-Boolean-Functions-by-Ryan-ODonnell}{
  pdf-ffbe9c7b3306-p019-b003,
  pdf-ffbe9c7b3306-p143-b004,
  pdf-ffbe9c7b3306-p143-b005
}
\label{FABL.ZeroOneFunction}
\leanok
For a natural number $n$, a zero-one function on the sign cube of length $n$ is a
function from that cube to $\{\text{true},\text{false}\}$.
\end{definition}

\begin{definition}[Real encoding of a Boolean value]
\lean{FABL.zeroOneValue}
\statementsource{Analysis-of-Boolean-Functions-by-Ryan-ODonnell}{
  pdf-ffbe9c7b3306-p143-b004,
  pdf-ffbe9c7b3306-p143-b005
}
\label{FABL.zeroOneValue}
\leanok
Encode a Boolean value as a real number by assigning $1$ to true and $0$ to false.
\end{definition}

\begin{definition}[Real-valued form of a Boolean function]
\lean{FABL.ZeroOneFunction.toReal}
\label{FABL.ZeroOneFunction.toReal}
\leanok
\uses{FABL.ZeroOneFunction, FABL.zeroOneValue}
For a function $f:\{−1,1\}^n\to\{\mathrm{false},\mathrm{true}\}$, its real-valued form
is the function on the same domain obtained by encoding each Boolean output as a real
number.
\end{definition}

\begin{definition}[Probability distribution of a biased bit]
\lean{FABL.pBiasedBitPMF}
\statementsource{Analysis-of-Boolean-Functions-by-Ryan-ODonnell}{
  pdf-ffbe9c7b3306-p017-b007,
  pdf-ffbe9c7b3306-p143-b005,
  pdf-ffbe9c7b3306-p220-b003
}
\label{FABL.pBiasedBitPMF}
\leanok
For a real number $p$ with $0\le p\le 1$, define a probability distribution on the
Boolean values that assigns probability $p$ to true and probability $1-p$ to false.
\end{definition}

\begin{theorem}[Mean of a biased bit]
\lean{FABL.pmfExpectation_pBiasedBitPMF_zeroOneValue}
\label{FABL.pmfExpectation_pBiasedBitPMF_zeroOneValue}
\leanok
\uses{FABL.pBiasedBitPMF, FABL.pmfExpectation, FABL.zeroOneValue}
\difficulty{3}
Let $0\leq p\leq 1$, and give the two-point set $\{\mathrm{false},\mathrm{true}\}$ the
probability law assigning mass $1-p$ to $\mathrm{false}$ and mass $p$ to
$\mathrm{true}$. The expectation of the indicator that the bit is true is $p$.
\end{theorem}

\begin{definition}[Biased random Boolean function]
\lean{FABL.pBiasedRandomFunctionPMF}
\statementsource{Analysis-of-Boolean-Functions-by-Ryan-ODonnell}{pdf-ffbe9c7b3306-p143-b005}
\label{FABL.pBiasedRandomFunctionPMF}
\leanok
\uses{FABL.ZeroOneFunction, FABL.independentProductPMF, FABL.pBiasedBitPMF}
For a natural number $n$ and $p\in[0,1]$, define a probability distribution on
Boolean-valued functions on the sign cube of length $n$ by choosing the output at each
input independently, with probability $p$ of true and probability $1-p$ of false.
\end{definition}

\begin{theorem}[Point probabilities of a biased random function]
\lean{FABL.pBiasedRandomFunctionPMF_apply}
\label{FABL.pBiasedRandomFunctionPMF_apply}
\leanok
\uses{FABL.ZeroOneFunction, FABL.pBiasedBitPMF, FABL.pBiasedRandomFunctionPMF}
\difficulty{1}
Let $n$ be a natural number, let $p\in[0,1]$, and let
$f:\{-1,1\}^n\to\{\mathrm{false},\mathrm{true}\}$. Under the product probability law in
which the values at distinct inputs are independent and each is true with probability
$p$, the probability of obtaining exactly $f$ is the product, over all inputs
$x\in\{-1,1\}^n$, of the probability that a single bit equals $f(x)$.
\end{theorem}

\begin{theorem}[Measure of a biased random function]
\lean{FABL.pBiasedRandomFunctionPMF_toMeasure}
\label{FABL.pBiasedRandomFunctionPMF_toMeasure}
\leanok
\uses{FABL.independentProductPMF_toMeasure, FABL.pBiasedBitPMF, FABL.pBiasedRandomFunctionPMF}
\difficulty{1}
Let $n$ be a natural number and let $p\in[0,1]$. The measure associated with the product
probability law on Boolean-valued functions on $\{-1,1\}^n$, with independent values
each true with probability $p$, equals the product measure whose coordinate laws are the
corresponding single-bit probability law.
\end{theorem}

\begin{theorem}[Mean of a biased bit]
\lean{FABL.integral_zeroOneValue_pBiasedBitPMF}
\label{FABL.integral_zeroOneValue_pBiasedBitPMF}
\leanok
\uses{FABL.pBiasedBitPMF, FABL.pmfExpectation_eq_integral, FABL.pmfExpectation_pBiasedBitPMF_zeroOneValue, FABL.zeroOneValue}
\difficulty{3}
Let $p\in[0,1]$, and let $b$ be a Boolean random variable that equals true with
probability $p$ and false with probability $1-p$. Assign value $1$ to true and $0$ to
false. The mean of this value is $p$.
\end{theorem}

\begin{theorem}[Square Integrability of a Biased Bit]
\lean{FABL.zeroOneValue_memLp_pBiasedBitPMF}
\label{FABL.zeroOneValue_memLp_pBiasedBitPMF}
\leanok
\uses{FABL.pBiasedBitPMF, FABL.zeroOneValue}
\difficulty{4}
For every real number $p$ with $0\le p\le 1$, the function on a Boolean bit that takes
value $1$ on true and $0$ on false belongs to $L^2$ with respect to the probability
measure assigning mass $p$ to true and $1-p$ to false.
\end{theorem}

\begin{theorem}[Variance of a biased bit]
\lean{FABL.variance_zeroOneValue_pBiasedBitPMF}
\label{FABL.variance_zeroOneValue_pBiasedBitPMF}
\leanok
\uses{FABL.integral_zeroOneValue_pBiasedBitPMF, FABL.pBiasedBitPMF, FABL.zeroOneValue, FABL.zeroOneValue_memLp_pBiasedBitPMF}
\difficulty{4}
For every $p\in[0,1]$, the variance of the zero-one value of a bit that equals $1$ with
probability $p$ and $0$ with probability $1-p$ is $p(1-p)$.
\end{theorem}

\begin{theorem}[Measure abs p biased weighted sum sub mean ge le]
\lean{FABL.measure_abs_pBiasedWeightedSum_sub_mean_ge_le}
\statementsource{Analysis-of-Boolean-Functions-by-Ryan-ODonnell}{
  pdf-ffbe9c7b3306-p143-b005,
  pdf-ffbe9c7b3306-p144-b001
}
\label{FABL.measure_abs_pBiasedWeightedSum_sub_mean_ge_le}
\leanok
\uses{FABL.ZeroOneFunction, FABL.integral_zeroOneValue_pBiasedBitPMF, FABL.pBiasedBitPMF, FABL.pBiasedRandomFunctionPMF, FABL.pBiasedRandomFunctionPMF_toMeasure, FABL.zeroOneValue}
\difficulty{6}
% dataset-pipeline: stub -- prose pending build_informal_statements.py
Statement pending regeneration.
\end{theorem}

\begin{theorem}[Centered weighted sum formula]
\lean{FABL.fourierCoeff_sub_pBiasedMean_eq_weightedCenteredSum}
\label{FABL.fourierCoeff_sub_pBiasedMean_eq_weightedCenteredSum}
\leanok
\uses{FABL.ZeroOneFunction, FABL.ZeroOneFunction.toReal, FABL.expect_monomial, FABL.fourierCoeff, FABL.monomial, FABL.zeroOneValue}
\difficulty{4}
Let $n$ be a natural number, let $p\in\mathbb R$, let $S$ be a subset of the $n$
coordinates, and let $f:\{-1,1\}^n\to\{0,1\}$. Its Fourier coefficient, computed
using the real-valued version of $f$, minus $p$ when $S$ is empty and minus zero
otherwise, equals the sum over all cube points $x$ of $\frac{\prod_{i\in
S}x_i}{2^n}(f(x)-p)$.
\end{theorem}

\begin{theorem}[Squared monomial sum]
\lean{FABL.sum_sq_normalized_monomial}
\label{FABL.sum_sq_normalized_monomial}
\leanok
\uses{FABL.monomial, FABL.monomial_sq}
\difficulty{3}
Let $n$ be a natural number and let $S$ be a subset of the $n$ coordinates. On the sign
cube, the sum over all points $x$ of the squared monomial at $S$, divided by the square
of the cube's cardinality, equals the reciprocal of the cube's cardinality.
\end{theorem}

\begin{theorem}[Variance of a random Fourier coefficient]
\lean{FABL.variance_pBiasedRandomFunction_fourierCoeff}
\label{FABL.variance_pBiasedRandomFunction_fourierCoeff}
\leanok
\uses{FABL.Sign, FABL.ZeroOneFunction, FABL.ZeroOneFunction.toReal, FABL.fourierCoeff, FABL.fourierCoeff_sub_pBiasedMean_eq_weightedCenteredSum, FABL.monomial, FABL.pBiasedBitPMF, FABL.pBiasedRandomFunctionPMF, FABL.pBiasedRandomFunctionPMF_toMeasure, FABL.sum_sq_normalized_monomial, FABL.variance_zeroOneValue_pBiasedBitPMF, FABL.zeroOneValue}
\difficulty{5}
Let $n$ be a natural number, let $p\in[0,1]$, and let $S\subseteq\{0,\ldots,n-1\}$ be
any set of coordinates. Choose a random Boolean function on the sign cube by assigning
independently to each input a value equal to $1$ with probability $p$ and $0$ with
probability $1-p$. The variance, over this random choice, of its Fourier coefficient
indexed by $S$ is $\frac{p(1-p)}{2^n}$.
\end{theorem}

\begin{theorem}[Concentration of a biased Fourier coefficient]
\lean{FABL.measure_pBiasedRandomFunction_fourierCoeff_sub_mean_ge_le}
\proofsource{Analysis-of-Boolean-Functions-by-Ryan-ODonnell}{
  pdf-ffbe9c7b3306-p143-b005,
  pdf-ffbe9c7b3306-p144-b001
}
\label{FABL.measure_pBiasedRandomFunction_fourierCoeff_sub_mean_ge_le}
\leanok
\uses{FABL.ZeroOneFunction.toReal, FABL.fourierCoeff, FABL.fourierCoeff_sub_pBiasedMean_eq_weightedCenteredSum, FABL.measure_abs_pBiasedWeightedSum_sub_mean_ge_le, FABL.monomial, FABL.pBiasedRandomFunctionPMF, FABL.sum_sq_normalized_monomial}
\difficulty{3}
Let $n\in\mathbb{N}$, let $p\in[0,1]$, let $S\subseteq[n]$, and let $t\ge0$. Choose a
random function $f:\{-1,1\}^n\to\{0,1\}$ by assigning independent values, each equal to
$1$ with probability $p$ and $0$ otherwise. The probability that the absolute difference
between its $S$-th Fourier coefficient and $p$ when $S=\varnothing$, or $0$ otherwise,
is at least $t$ is at most $2\exp\!\left(-2^n t^2/2\right)$.
\end{theorem}

\begin{theorem}[Tail bound for a biased Fourier coefficient]
\lean{FABL.measure_pBiasedRandomFunction_fourierCoeff_atThreshold_ge_le}
\proofsource{Analysis-of-Boolean-Functions-by-Ryan-ODonnell}{
  pdf-ffbe9c7b3306-p143-b005,
  pdf-ffbe9c7b3306-p144-b001
}
\label{FABL.measure_pBiasedRandomFunction_fourierCoeff_atThreshold_ge_le}
\leanok
\uses{FABL.ZeroOneFunction.toReal, FABL.card_mul_randomBooleanFourierThreshold_sq_div_two, FABL.fourierCoeff, FABL.measure_pBiasedRandomFunction_fourierCoeff_sub_mean_ge_le, FABL.pBiasedRandomFunctionPMF, FABL.randomBooleanFourierThreshold, FABL.randomBooleanFourierThreshold_nonneg}
\difficulty{3}
Let $n\in\mathbb{N}$, let $p\in[0,1]$, and let $S\subseteq[n]$. Choose a random function
$f:\{-1,1\}^n\to\{0,1\}$ by assigning independent values, each equal to $1$ with
probability $p$ and $0$ otherwise. The probability that the absolute difference between
its $S$-th Fourier coefficient (after viewing its values as real numbers) and $p$ when
$S=\varnothing$, or $0$ otherwise, is at least $2\sqrt n\,2^{-n/2}$ is at most
$2\exp(-2n)$.
\end{theorem}

\begin{theorem}[Containment of Fourier failure sets]
\lean{FABL.pBiasedFourierFailure_subset}
\label{FABL.pBiasedFourierFailure_subset}
\leanok
\uses{FABL.ZeroOneFunction, FABL.ZeroOneFunction.toReal, FABL.fourierCoeff, FABL.randomBooleanFourierThreshold}
\difficulty{4}
Let $n$ be a natural number, let $p$ be real, and let $\tau_n=2\sqrt n\,2^{-n/2}$. For a
Boolean-valued function $f$ on the $n$-dimensional Boolean cube, write $f_{\mathbb R}$
for its real-valued indicator and $\widehat f_{\mathbb R}(S)$ for its Fourier
coefficient at $S$. If either $|\widehat f_{\mathbb R}(\varnothing)-p|>\tau_n$ or some
nonempty coordinate set $S$ satisfies $|\widehat f_{\mathbb R}(S)|>\tau_n$, then there
is a coordinate set $S$ such that
\[
\left|\widehat f_{\mathbb
R}(S)-\begin{cases}p,&S=\varnothing,\\0,&S\ne\varnothing\end{cases}\right|\ge\tau_n.
\]
\end{theorem}

\begin{theorem}[Simultaneous Fourier failure probability]
\lean{FABL.measure_pBiasedRandomFunction_fourierFailure_le}
\proofsource{Analysis-of-Boolean-Functions-by-Ryan-ODonnell}{
  pdf-ffbe9c7b3306-p143-b005,
  pdf-ffbe9c7b3306-p144-b001
}
\label{FABL.measure_pBiasedRandomFunction_fourierFailure_le}
\leanok
\uses{FABL.ZeroOneFunction, FABL.ZeroOneFunction.toReal, FABL.fourierCoeff, FABL.measure_pBiasedRandomFunction_fourierCoeff_atThreshold_ge_le, FABL.pBiasedFourierFailure_subset, FABL.pBiasedRandomFunctionPMF, FABL.randomBooleanFourierThreshold, FABL.two_pow_mul_two_exp_neg_two_mul_le_rpow_neg}
\difficulty{4}
Let $n\in\mathbb{N}$ with $1<n$, and let $p\in[0,1]$. Choose a random function
$f:\{-1,1\}^n\to\{0,1\}$ by assigning independent values, each equal to $1$ with
probability $p$ and $0$ otherwise. Put $\tau_n=2\sqrt n\,2^{-n/2}$. The probability that
either the constant Fourier coefficient differs from $p$ by more than $\tau_n$, or some
nonempty Fourier index has coefficient of absolute value greater than $\tau_n$, is at
most $2^{-n}$.
\end{theorem}

\begin{theorem}[Fourier regularity of a random Boolean function]
\lean{FABL.measure_pBiasedRandomFunction_not_isFourierRegular_le}
\proofsource{Analysis-of-Boolean-Functions-by-Ryan-ODonnell}{
  pdf-ffbe9c7b3306-p143-b005,
  pdf-ffbe9c7b3306-p144-b001,
  pdf-ffbe9c7b3306-p144-b003
}
\label{FABL.measure_pBiasedRandomFunction_not_isFourierRegular_le}
\leanok
\uses{FABL.IsFourierRegular, FABL.ZeroOneFunction.toReal, FABL.measure_pBiasedRandomFunction_fourierFailure_le, FABL.pBiasedRandomFunctionPMF, FABL.randomBooleanFourierThreshold}
\difficulty{4}
Let $n$ be an integer with $1<n$, and let $p\in[0,1]$. Choose a Boolean function on the
$n$-dimensional sign cube by assigning each input an independent Bernoulli-$p$ bit, and
regard its values as real numbers $0$ and $1$. The probability that its nonempty Fourier
coefficients do not all have absolute value at most $2\sqrt{n}\,2^{-n/2}$ is at most
$2^{-n}$.
\end{theorem}

\begin{definition}[Boolean parity sign]
\lean{FABL.randomBooleanParitySign}
\label{FABL.randomBooleanParitySign}
\leanok
\uses{FABL.Sign}
For a subset $S\subseteq\{0,\ldots,n-1\}$ and $x\in\{−1,1\}^n$, define the parity sign
on $S$ to be $\prod_{i\in S}x_i$, regarded as a unit of the integers.
\end{definition}

\begin{theorem}[Value of a random Boolean parity sign]
\lean{FABL.signValue_randomBooleanParitySign}
\label{FABL.signValue_randomBooleanParitySign}
\leanok
\uses{FABL.monomial, FABL.randomBooleanParitySign, FABL.signValue}
\difficulty{1}
Let $n$ be a natural number, let $S\subseteq\{1,\ldots,n\}$, and let $x\in\{-1,1\}^n$.
The real value of the sign obtained by multiplying the coordinates $x_i$ for $i\in S$
equals the monomial $\prod_{i\in S}x_i$.
\end{theorem}

\begin{definition}[Boolean parity modulation]
\lean{FABL.randomBooleanModulationEquiv}
\label{FABL.randomBooleanModulationEquiv}
\leanok
\uses{FABL.BooleanFunction, FABL.randomBooleanParitySign}
For a subset $S\subseteq\{0,\ldots,n-1\}$, define a bijection on functions
$f:\{−1,1\}^n\to\{−1,1\}$ by sending $f$ to the function $x\mapsto f(x)\prod_{i\in
S}x_i$.
\end{definition}

\begin{theorem}[Modulation and Fourier coefficients]
\lean{FABL.fourierCoeff_randomBooleanModulation_empty}
\label{FABL.fourierCoeff_randomBooleanModulation_empty}
\leanok
\uses{FABL.BooleanFunction, FABL.BooleanFunction.toReal, FABL.fourierCoeff, FABL.monomial, FABL.randomBooleanModulationEquiv, FABL.randomBooleanParitySign, FABL.signValue, FABL.signValue_randomBooleanParitySign}
\difficulty{3}
Let $n$ be a nonnegative integer, let $S$ be a subset of the $n$ coordinates, and let
$f$ be a function from the sign cube $\{−1,1\}^n$ to $\{−1,1\}$. Multiplying $f$
pointwise by the parity character on $S$ makes its Fourier coefficient at the empty set
equal to the Fourier coefficient of $f$ at $S$.
\end{theorem}

\begin{theorem}[Second moment of a random Fourier coefficient]
\lean{FABL.expect_sq_fourierCoeff_uniformBooleanFunction}
\proofsource{Analysis-of-Boolean-Functions-by-Ryan-ODonnell}{pdf-ffbe9c7b3306-p035-b007}
\label{FABL.expect_sq_fourierCoeff_uniformBooleanFunction}
\leanok
\uses{FABL.BooleanFunction, FABL.BooleanFunction.toReal, FABL.fourierCoeff, FABL.fourierCoeff_randomBooleanModulation_empty, FABL.randomBooleanModulationEquiv, FABL.sum_sq_fourierCoeff_eq_one}
\difficulty{5}
Let $n$ be a natural number and let $S$ be any subset of the $n$ coordinates. Choose a
Boolean function from the $n$-dimensional sign cube to the signs uniformly from all such
functions, and regard its values as real numbers. The expected square of its Fourier
coefficient at $S$ is $2^{-n}$.
\end{theorem}

\begin{definition}[Bijection by deleting a specified element]
\lean{FABL.randomFunctionEraseContainingEquiv}
\label{FABL.randomFunctionEraseContainingEquiv}
\leanok
Let $n$ be a nonnegative integer and let $i$ be an index with $0\le i<n$. Deleting $i$
defines a bijection from the subsets of $\{0,\ldots,n-1\}$ that contain $i$ to those
that do not contain $i$; its inverse adjoins $i$.
\end{definition}

\begin{theorem}[Weighted sum over subsets containing a coordinate]
\lean{FABL.sum_pow_card_sub_one_filter_mem}
\label{FABL.sum_pow_card_sub_one_filter_mem}
\leanok
\uses{FABL.randomFunctionEraseContainingEquiv}
\difficulty{5}
Let $n$ be a natural number, let $\rho\in\mathbb R$, and let $i$ be one of the $n$
coordinates. The sum of $\rho^{|S|-1}$ over all subsets $S$ of the coordinates that
contain $i$ equals $(1+\rho)^{n-1}$.
\end{theorem}

\begin{theorem}[Mean stable influence of a uniform Boolean function]
\lean{FABL.pmfExpectation_uniformBooleanFunction_stableInfluence}
\proofsource{Analysis-of-Boolean-Functions-by-Ryan-ODonnell}{pdf-ffbe9c7b3306-p145-b006}
\label{FABL.pmfExpectation_uniformBooleanFunction_stableInfluence}
\leanok
\uses{FABL.BooleanFunction, FABL.BooleanFunction.toReal, FABL.expect_sq_fourierCoeff_uniformBooleanFunction, FABL.pmfExpectation, FABL.pmfExpectation_uniformPMF_eq_expect, FABL.stableInfluence, FABL.sum_pow_card_sub_one_filter_mem, FABL.uniformPMF}
\difficulty{3}
Let $n\ge 1$, let $\rho\in\mathbb{R}$, and choose a coordinate $i\in\{1,\ldots,n\}$. A
Boolean function assigns a sign in $\{-1,1\}$ to each sign vector in $\{-1,1\}^n$. The
average, over all such functions with the uniform distribution, of the stable influence
at coordinate $i$ is $\frac{(1+\rho)^{n-1}}{2^n}$.
\end{theorem}

\begin{theorem}[Mean stable influence of a random Boolean function]
\lean{FABL.pmfExpectation_uniformBooleanFunction_stableInfluence_one_sub}
\proofsource{Analysis-of-Boolean-Functions-by-Ryan-ODonnell}{pdf-ffbe9c7b3306-p145-b006}
\label{FABL.pmfExpectation_uniformBooleanFunction_stableInfluence_one_sub}
\leanok
\uses{FABL.BooleanFunction, FABL.BooleanFunction.toReal, FABL.pmfExpectation, FABL.pmfExpectation_uniformBooleanFunction_stableInfluence, FABL.stableInfluence, FABL.uniformPMF}
\difficulty{3}
Let $n$ be a nonnegative integer, let $\delta\in[0,1]$, and choose uniformly a function
from the Boolean cube $\{-1,1\}^n$ to $\{\pm1\}$. For each coordinate
$i\in\{1,\ldots,n\}$, the expected stable influence of this function at noise parameter
$1-\delta$ is $\frac{(1-\delta/2)^n}{2-\delta}$.
\end{theorem}
